% ============================================================
%  Chapter 16 — Determinant
% ============================================================
\section{Determinant}
\label{sec:determinant}

The determinant is a scalar-valued function on square matrices (or
square-matrix-valued linear operators) that encodes essential geometric and
algebraic information: it measures the signed volume scaling of a linear map
and vanishes precisely when the matrix is singular.

% ---- Motivating example: 2×2 invertibility ------------------
\begin{tipbox}[Remark --- Motivation: Invertibility of $2\times 2$ Matrices]
A matrix $A=\begin{pmatrix}a&b\\c&d\end{pmatrix}$ is invertible if and
only if $ad-bc\neq 0$.  This can be verified by row-reducing the augmented
matrix $[A\mid I_2]$.  The scalar $ad-bc$ is precisely the
\emph{determinant} of $A$.
\end{tipbox}

% ---- Multilinear and alternating ----------------------------
\begin{defbox}[Definition --- Multilinear and Alternating Function]
\label{def:multilinear-alternating}
Let $A=(A_1\mid A_2\mid\cdots\mid A_n)\in\F^{n\times n}$ be written in
column form.  A function $f:\F^{n\times n}\to\F$ is called:
\begin{enumerate}[label=\textbf{[\roman*]}]
    \item \textbf{Multilinear (in columns)} if, for every $1\le i\le n$
          and every $\alpha\in\F$ and columns $A_i,A'_i\in\F^n$:
          \[
            f(A_1\mid\cdots\mid\alpha A_i+A'_i\mid\cdots\mid A_n)
            = \alpha\,f(A_1\mid\cdots\mid A_i\mid\cdots\mid A_n)
            + f(A_1\mid\cdots\mid A'_i\mid\cdots\mid A_n).
          \]
    \item \textbf{Alternating (in columns)} if
          $f(\ldots\mid A_i\mid\cdots\mid A_j\mid\ldots)=0$
          whenever $A_i=A_j$ for some $i\neq j$.
\end{enumerate}
A function that is both alternating and multilinear is called an
\textbf{alternating multilinear function} (or a \textbf{form}).
If additionally $f(I_n)\neq 0$, it is called a \textbf{volume form}.
\end{defbox}

% ---- Definition of determinant ------------------------------
\begin{defbox}[Definition --- Determinant]
\label{def:determinant}
The unique alternating multilinear function
$\det:\F^{n\times n}\to\F$ satisfying $\det(I_n)=1$ is called the
\textbf{determinant}.  Alternative notation: $\det(A)=|A|$.
\end{defbox}

% ---- Determinant of a linear operator -----------------------
\begin{defbox}[Definition --- Determinant of a Linear Operator]
\label{def:det-lt}
Let $f:V\to V$ be a linear operator on a finite-dimensional vector space
$V$ over $\F$ and let $\mathcal{B}$ be any ordered basis for $V$.
Define
\[
    \det(f) \;=\; \det\bigl([f]_{\mathcal{B}}\bigr).
\]
\end{defbox}

\begin{tipbox}[Remark --- Determinant is Basis-Independent]
The determinant of a linear operator is the same for every choice of
basis.  Indeed, if $\mathcal{B}$ and $\mathcal{B}'$ are two ordered bases,
then $[f]_{\mathcal{B}'}=C^{-1}[f]_{\mathcal{B}}C$ for some invertible $C$,
and $\det(C^{-1}[f]_{\mathcal{B}}C)=\det([f]_{\mathcal{B}})$ (see the
multiplicativity property below).
\end{tipbox}

% ---- Characterisation of alt. mult. functions ---------------
\begin{propbox}[Proposition --- Characterisation of Alternating Multilinear Functions]
\label{prop:alt-mult-det}
Let $f:\F^{n\times n}\to\F$ be an alternating multilinear function with
$f(I_n)=\alpha$.  Then
\[
    f = \alpha\cdot\det.
\]
\end{propbox}

\begin{corbox}[Corollary --- Existence of the Determinant]
For each $n\in\N$, there exists an alternating multilinear function
$\det:\F^{n\times n}\to\F$ (i.e.\ the determinant exists).
\end{corbox}

\begin{thmbox}[Theorem --- Uniqueness of Alternating Multilinear Functions]
\label{thm:unique-alt-mult}
For each $n\in\N$ and each $\alpha\in\F$, there exists a \emph{unique}
alternating multilinear function $f:\F^{n\times n}\to\F$ such that
$f(I_n)=\alpha$.
\end{thmbox}

\begin{corbox}[Corollary --- Uniqueness of the Determinant]
For each $n\in\N$, the determinant is the unique alternating multilinear
function satisfying $f(I_n)=1$.
\end{corbox}

% ---- 2×2 determinant ----------------------------------------
\begin{exbox}[Example --- Determinant of a $2\times 2$ Matrix]
\label{ex:det-2x2}
The function $f:\F^{2\times 2}\to\F$ defined by
\[
    f\!\begin{pmatrix}a&b\\c&d\end{pmatrix} = ad - bc
\]
is alternating and multilinear.  Since $f(I_2)=1\cdot 1-0\cdot 0=1$, it
is the determinant on $\F^{2\times 2}$.

\smallskip\noindent
The three families of alternating multilinear functions on
$\F^{2\times 2}$ illustrate the characterisation above:
\begin{enumerate}
    \item $n=1$: $f:\F^{1\times 1}\to\F$ is alternating multilinear if
          and only if $f$ is a linear map.
    \item $n=2$: $f\begin{pmatrix}a&b\\c&d\end{pmatrix}=ad-bc$
          (the determinant, $\alpha=1$).
    \item $n=2$: $f\begin{pmatrix}a&b\\c&d\end{pmatrix}=\alpha(ad-bc)$
          for any $\alpha\in\F$ (general form with $f(I_2)=\alpha$).
\end{enumerate}
\end{exbox}

% ---- Properties of alternating multilinear functions --------
\begin{propbox}[Proposition --- Properties of Alternating Multilinear Functions]
\label{prop:alt-mult-props}
Let $f:\F^{n\times n}\to\F$ be an alternating multilinear function.
Then:
\begin{enumerate}[label=\textbf{(\alph*)}]
    \item Swapping two columns changes the sign:
          \[
            f(\ldots\mid A_i\mid\cdots\mid A_j\mid\ldots)
            = -f(\ldots\mid A_j\mid\cdots\mid A_i\mid\ldots).
          \]
    \item A zero column gives zero:
          \[
            f(A_1\mid\cdots\mid 0\mid\cdots\mid A_n)=0.
          \]
    \item Adding a scalar multiple of one column to another is neutral:
          \[
            f(\ldots\mid A_i+\alpha A_j\mid\cdots\mid A_j\mid\ldots)
            = f(\ldots\mid A_i\mid\cdots\mid A_j\mid\ldots)
            \quad\text{for all }\alpha\in\F,\; i\neq j.
          \]
    \item If $A_i=\sum_{j\neq i}\alpha_j A_j$ (column is a linear
          combination of the others), then $f(A)=0$.
\end{enumerate}
\end{propbox}

\begin{corbox}[Corollary --- Rank Condition for Alternating Multilinear Functions]
If $f:\F^{n\times n}\to\F$ is alternating multilinear and
$\operatorname{rank}A < n$, then $f(A)=0$.
\end{corbox}

% ---- Submatrix notation -------------------------------------
\begin{defbox}[Definition --- Submatrix $A(i\mid j)$]
\label{def:submatrix}
Given $A\in\F^{n\times n}$ and indices $1\le i,j\le n$, the
\textbf{$(i,j)$-minor matrix} $A(i\mid j)\in\F^{(n-1)\times(n-1)}$
is the matrix obtained from $A$ by deleting row $i$ and column $j$.
\end{defbox}

% ---- Lemma on dimension reduction ---------------------------
\begin{thmbox}[Lemma --- Constructing Alt.\ Mult.\ Functions by Dimension Reduction]
\label{lem:dim-reduction}
Let $f:\F^{(n+1)\times(n+1)}\to\F$ be alternating multilinear with
$f(I_{n+1})=\alpha$.  For $i=1,\ldots,n$ define
$f_i:\F^{n\times n}\to\F$ by $f_i(A)=f(A_{(i)})$, where $A_{(i)}$ is
the $(n+1)\times(n+1)$ matrix
\[
    A_{(i)} =
    \begin{pmatrix}
        a_{1,1}&\cdots&a_{1,i-1}&0&a_{1,i}&\cdots&a_{1,n}\\
        \vdots & & \vdots & \vdots & \vdots & & \vdots \\
        a_{i-1,1}&\cdots&a_{i-1,i-1}&0&a_{i-1,i}&\cdots&a_{i-1,n}\\
        0&\cdots&0&1&0&\cdots&0\\
        a_{i,1}&\cdots&a_{i,i-1}&0&a_{i,i}&\cdots&a_{i,n}\\
        \vdots & & \vdots & \vdots & \vdots & & \vdots \\
        a_{n,1}&\cdots&a_{n,i-1}&0&a_{n,i}&\cdots&a_{n,n}
    \end{pmatrix},
    \quad A=(a_{ij})\in\F^{n\times n}.
\]
Then each $f_i$ is an alternating multilinear function from
$\F^{n\times n}$ to $\F$.
\end{thmbox}

% ---- Expansion formulas -------------------------------------
\begin{flowbox}[Key Formula --- Determinant Expansion by Row or Column]
\label{flow:det-expansion}
For $A\in\F^{n\times n}$ and any fixed row $i$ or column $j$:
\begin{align*}
    \det(A) &= \sum_{j=1}^{n}(-1)^{i+j}\,a_{ij}\,\det\!\bigl(A(i\mid j)\bigr)
               &&\text{(expansion along row }i\text{)}\\[4pt]
    \det(A) &= \sum_{i=1}^{n}(-1)^{i+j}\,a_{ij}\,\det\!\bigl(A(i\mid j)\bigr)
               &&\text{(expansion along column }j\text{)}
\end{align*}
The number $(-1)^{i+j}\det(A(i\mid j))$ is called the \textbf{$(i,j)$-cofactor} of $A$.

\medskip
\textbf{Explicit formulas:}
\[
    \begin{vmatrix}a_{11}&a_{12}\\a_{21}&a_{22}\end{vmatrix}
    = a_{11}a_{22}-a_{12}a_{21}.
\]
\begin{align*}
    \begin{vmatrix}
        a_{11}&a_{12}&a_{13}\\a_{21}&a_{22}&a_{23}\\a_{31}&a_{32}&a_{33}
    \end{vmatrix}
    &= a_{11}\begin{vmatrix}a_{22}&a_{23}\\a_{32}&a_{33}\end{vmatrix}
      -a_{12}\begin{vmatrix}a_{21}&a_{23}\\a_{31}&a_{33}\end{vmatrix}
      +a_{13}\begin{vmatrix}a_{21}&a_{22}\\a_{31}&a_{32}\end{vmatrix}
      &&(\text{row }1)\\[4pt]
    &= a_{11}\begin{vmatrix}a_{22}&a_{23}\\a_{32}&a_{33}\end{vmatrix}
      -a_{21}\begin{vmatrix}a_{12}&a_{13}\\a_{32}&a_{33}\end{vmatrix}
      +a_{31}\begin{vmatrix}a_{12}&a_{13}\\a_{22}&a_{23}\end{vmatrix}
      &&(\text{col }1)
\end{align*}
\end{flowbox}

\begin{propbox}[Proposition --- Constructing Det from Lower-Dimensional Det]
\label{prop:det-from-lower}
Let $g:\F^{(n-1)\times(n-1)}\to\F$ be an alternating multilinear function
with $g(I_{n-1})=\alpha$, and let $n\ge 2$.  Then
\[
    f(A) = \sum_{j=1}^{n}(-1)^{i+j}\,a_{ij}\,g\!\bigl(A(i\mid j)\bigr)
\]
(for any fixed $i$) defines an alternating multilinear function
$f:\F^{n\times n}\to\F$ with $f(I_n)=\alpha$.
\end{propbox}

\begin{methodbox}[Method --- Efficient Determinant Calculation]
When computing $\det(A)$ by cofactor expansion, choose the row or column
that contains the most zeros to minimise the number of terms.
\end{methodbox}

% ---- Properties of the determinant --------------------------
\begin{propbox}[Proposition --- Determinant of the Transpose]
\label{prop:det-transpose}
For all $A\in\F^{n\times n}$:
\[
    \det(A) = \det(A^T).
\]
\end{propbox}

\begin{propbox}[Proposition --- Effect of Elementary Row Operations on the Determinant]
\label{prop:det-row-ops}
Let $A\in\F^{n\times n}$ and let $A'$ be obtained from $A$ by an
elementary row operation.
\begin{enumerate}[label=\textbf{(\roman*)}]
    \item \emph{Type~1} ($r_i\mapsto\alpha r_i$, $\alpha\neq 0$):
          $\det(A')=\alpha\det(A)$.
    \item \emph{Type~2} ($r_i\mapsto r_i+\alpha r_j$, $i\neq j$):
          $\det(A')=\det(A)$.
    \item \emph{Type~3} ($r_i\leftrightarrow r_j$, $i\neq j$):
          $\det(A')=-\det(A)$.
\end{enumerate}
\end{propbox}

\begin{propbox}[Proposition --- Determinant of an Upper-Triangular Matrix]
\label{prop:det-triangular}
If $A=(a_{ij})\in\F^{n\times n}$ is upper triangular (i.e.\
$a_{ij}=0$ for $i>j$), then
\[
    \det(A) = \prod_{i=1}^{n} a_{ii}.
\]
\end{propbox}

\begin{corbox}[Corollary --- Determinant of a Diagonal Matrix]
If $D=\operatorname{diag}(a_{11},a_{22},\ldots,a_{nn})$, then
$\det(D)=\prod_{i=1}^n a_{ii}$.
\end{corbox}

\begin{propbox}[Proposition --- Multiplicativity of the Determinant]
\label{prop:det-product}
For $A,B\in\F^{n\times n}$:
\[
    \det(AB) = \det(A)\cdot\det(B).
\]
\end{propbox}

\begin{corbox}[Corollary --- Further Properties of the Determinant]
\label{cor:det-properties}
For $A,B\in\F^{n\times n}$ and $\alpha\in\F$:
\begin{enumerate}[label=\textbf{(\roman*)}]
    \item $\det(AB)=\det(BA)$.
    \item If $B$ is invertible: $\det(BAB^{-1})=\det(A)$.  In particular,
          \emph{similar matrices have the same determinant}.
    \item $\det(\alpha A)=\alpha^n\det(A)$.
\end{enumerate}
\end{corbox}

\begin{propbox}[Proposition --- Invertibility and the Determinant]
\label{prop:det-invertible}
$A\in\F^{n\times n}$ is invertible if and only if $\det(A)\neq 0$.
\end{propbox}

\begin{corbox}[Corollary --- Determinant of an Inverse Matrix]
If $A$ is invertible, then
\[
    \det(A^{-1}) = \frac{1}{\det(A)}.
\]
\end{corbox}

\begin{warnbox}[Warning --- Determinant is NOT Additive]
In general, $\det(A+B)\neq\det(A)+\det(B)$.

\smallskip\noindent
\textbf{Counterexample.} Take $A=I_2$ and $B=-I_2$.  Then
$A+B=O_2$, so $\det(A+B)=0$, but $\det(A)+\det(B)=1+1=2\neq 0$.
\end{warnbox}

% ---- Adjoint ------------------------------------------------
\subsection{The Adjoint Matrix}

\begin{defbox}[Definition --- Adjoint (Classical Adjoint) Matrix]
\label{def:adjoint}
Given $A\in\F^{n\times n}$, the \textbf{adjoint} (or \textbf{classical
adjoint}, or \textbf{adjugate}) of $A$ is the matrix
$\operatorname{adj}A\in\F^{n\times n}$ defined by
\[
    (\operatorname{adj}A)_{ij} = (-1)^{i+j}\det\!\bigl(A(j\mid i)\bigr).
\]
Note the transposition: the $(i,j)$-entry of $\operatorname{adj}A$ is the
$(j,i)$-cofactor of $A$.
\end{defbox}

\begin{propbox}[Proposition --- Adjoint--Inverse Relationship]
\label{prop:adjoint-inverse}
For any $A\in\F^{n\times n}$:
\[
    A\cdot\operatorname{adj}A
    = \operatorname{adj}A\cdot A
    = \det(A)\cdot I_n.
\]
Consequently, if $A$ is invertible, then
\[
    A^{-1} = \frac{1}{\det(A)}\operatorname{adj}A.
\]
\end{propbox}

\begin{exbox}[Example --- Computing the Adjoint and Inverse]
\label{ex:adjoint}
Let $A=\begin{pmatrix}1&5&2\\2&1&1\\1&2&0\end{pmatrix}$.

\medskip\noindent
\textbf{Computing one entry:}
\[
    (\operatorname{adj}A)_{13}
    = (-1)^{1+3}\det(A(3\mid 1))
    = (+1)\det\!\begin{pmatrix}5&2\\1&1\end{pmatrix}
    = 5-2 = 3.
\]

\medskip\noindent
\textbf{Full adjoint:}
\[
    \operatorname{adj}A =
    \begin{pmatrix}-2&4&3\\1&-2&3\\3&3&-9\end{pmatrix}.
\]

\medskip\noindent
\textbf{Verification:}
\[
    A\cdot\operatorname{adj}A
    = \begin{pmatrix}9&0&0\\0&9&0\\0&0&9\end{pmatrix}
    = 9\,I_3, \qquad \det(A)=9.
\]
Hence $A^{-1} = \dfrac{1}{9}\operatorname{adj}A$.
\end{exbox}
